The original paper evaluates the method on eight UCI datasets: Iris, AlcoholQCM, Seeds, Wine, Segment, ElectricalGrid, Avila, and Letter~\cite{zhang2021gmm}. Missing entries are generated synthetically from fully observed data at ratios from $10\%$ to $70\%$. The algorithm is compared against three static imputers, namely mean filling, zero filling, and EM filling, as well as dynamic K-means filling variants from earlier work~\cite{wang2019kmeans}.

\begin{table}[t]
    \caption{A stronger experimental protocol for a modern rewrite}
    \label{tab:protocol}
    \centering
    \begin{tabularx}{\linewidth}{@{}lX@{}}
        \toprule
        Item & Recommendation \\
        \midrule
        Missingness regimes & Evaluate MCAR, MAR, and stress-test MNAR-inspired corruption \\
        Cluster count & Report results with true $\scal{K}$ and model-selected $\scal{K}$ \\
        Metrics & ACC, NMI, F-score, purity, runtime, iterations, failure rate \\
        Stability & Repeat across multiple seeds and missing-pattern draws \\
        Ablations & Initialization, covariance type, regularization, stopping tolerance \\
        Baselines & Include both classical and modern incomplete-data clustering methods \\
        \bottomrule
    \end{tabularx}
\end{table}

\subsection{What the Reported Results Actually Show}
Across ACC, NMI, F-score, and purity, the paper consistently reports that the proposed method outperforms classical two-stage pipelines. The strongest pattern is not simply ``GMM is better than K-means.'' A more precise reading is:
\begin{itemize}
    \item when imputation is decoupled from clustering, the filled matrix can erase discriminative structure;
    \item when imputation is updated in tandem with cluster assignments, the filled values become more cluster-aware;
    \item the advantage grows as the missing ratio becomes moderate or high, where naive imputers collapse faster.
\end{itemize}

This interpretation is more informative than listing performance tables. It connects the empirical results back to the modeling decision introduced in the formulation.

\subsection{Local Reproduction on \texttt{CC GENERAL}}
Our project extends beyond the original paper by reproducing the method on a separate real tabular dataset, \texttt{CC GENERAL}, through the notebook \texttt{GMM\_Nhat.ipynb}. This dataset contains $8950$ credit-card customers and $17$ numeric behavioral features after dropping the customer identifier. In the stored notebook run, we inject approximately $30\%$ missingness, which yields $45{,}789$ missing entries.

The reproduction notebook then walks through the same logic as the paper:
\begin{itemize}
    \item initialize missing values by mean imputation;
    \item standardize the completed matrix;
    \item initialize GMM parameters with K-means;
    \item compute responsibilities and perform one explicit EM-style update;
    \item run the full alternating routine with \texttt{fit\_scale()}.
\end{itemize}

The recorded run converges after $199$ iterations and yields a silhouette score of $0.1489$. We do not claim that this value is intrinsically high; its importance is different. It demonstrates that the model can be trained end-to-end on a moderate-dimensional financial dataset with large injected missingness while still producing a coherent partition structure.

\begin{figure}[t]
    \centering
    \includegraphics[width=\linewidth]{notebook_exports/GMM_Nhat_cell8_out1.png}
    \caption{Reproduction artifact from \texttt{GMM\_Nhat.ipynb}: clustered samples after dynamic imputation, with previously incomplete rows highlighted. This figure is useful because it exposes where imputed samples land in the learned geometry.}
    \label{fig:cc-general-reproduction}
\end{figure}

\subsection{Why Dynamic Imputation Helps}
Suppose two clusters overlap in marginal feature means but differ in multivariate shape. A static imputer based on global averages or independent feature heuristics may pull incomplete samples toward a global center. In contrast, the one-stage GMM method uses posterior responsibilities and covariance structure to complete missing values in a way that remains compatible with latent cluster geometry. In that sense, the algorithm does not merely repair missingness; it \highlight{preserves separability}.

\subsection{Limits of the Experimental Design}
Even though the empirical section is convincing at first glance, several limitations should be acknowledged explicitly in a stronger report:
\begin{itemize}
    \item the missingness is synthetically generated from complete datasets, which primarily matches MCAR-style setups;
    \item the number of clusters is assumed known in advance, removing an important practical difficulty;
    \item there is no ablation on initialization, covariance regularization, or sensitivity to nearly singular clusters;
    \item the baselines do not include more recent incomplete-data clustering methods or modern probabilistic deep approaches;
    \item performance gains are shown, but runtime and memory tradeoffs are not analyzed with the same care.
\end{itemize}

\subsection{What a Better Reproduction Should Add}
A stronger reproduction should go beyond reproducing aggregate scores. It should also reveal \emph{why} the method wins and \emph{when} it fails. In practical terms, the following additions would materially improve the experimental section:
\begin{itemize}
    \item report per-dataset runtime and iteration counts to quantify the price paid for one-stage optimization;
    \item compare full covariance against diagonal and tied covariance variants to isolate the role of geometric flexibility;
    \item measure imputation error together with clustering quality to see whether better clustering always aligns with better reconstruction;
    \item examine sensitivity under class imbalance, since mixture models can behave differently when some components are weakly represented.
\end{itemize}

Our current notebooks already move in that direction because they preserve intermediate outputs, convergence summaries, and visual artifacts. The next step is to convert those artifacts into a systematic experimental log rather than leaving them as notebook-only evidence.

\subsection{A Curious Observation Worth Preserving}
The paper notes that for some datasets, performance with moderate missingness can exceed performance on the complete data. This is counterintuitive, but not impossible. A plausible explanation is that the iterative completion step can smooth noise or attenuate outliers by pulling uncertain coordinates toward cluster-consistent regions. That does not mean missingness is beneficial in general; rather, it suggests the algorithm may implicitly act as a form of regularized denoising. This is one of the most interesting observations in the article and deserves more attention than it receives in the original presentation.
